%
% ---------------------------------------------------------------------------------------
% VALOR T2K 5 sample joint fit analysis, Run 1-8
% ---------------------------------------------------------------------------------------
%
\documentclass[11pt]{elsarticle}

\usepackage{adjustbox}
\usepackage[toc,page]{appendix}       	% Appendix
\usepackage{amsfonts}                 	% Extra symbols
\usepackage{amsmath}
\usepackage{appendix}
\usepackage{booktabs}
\usepackage{caption}			% Use when minipaging figures & tables together
\usepackage{color}
\usepackage{dcolumn}					% align table columns on decimal point
\usepackage[figuresright]{rotating}		% if you have landscape tables
\usepackage{float}
\usepackage[margin=2.50cm]{geometry}  	% Tweak margins
\usepackage{graphicx}                 	% PostScript figures
\usepackage{hyperref} 			% Linkable contents and reference without colour
\usepackage{lineno}                   	% For printing the line numbers
\usepackage{listings}
\usepackage{longtable}
\usepackage{multirow}                 	% Fancy tables
\usepackage{natbib}                   	% For building bibliography with \cite{}
\usepackage{pdflscape}                   	% allow page to be landscape for large tables
\usepackage{rotate}                     % rotate text in tables
\usepackage{siunitx}					% Easier units exponentials etc
\usepackage{subcaption}					% Allows use of subfigure
\usepackage[table,dvipsnames]{xcolor}	% Fancy tables and colors
\usepackage{xspace}		        % Appropriate insertion of spaces after commands
\usepackage{cleveref}                   % Load after the hyperref package
\usepackage{pdflscape}
\usepackage{setspace}
\onehalfspacing
\DeclareSIUnit\POT{POT}

\newcommand{\nueCConepi}{\(\nue\) CC\( 1\pi^{+} \)}

% By default, LaTeX doesn't like to fill more than 0.7 of a text page with tables and graphics, 
% nor does it like too many figures per page. This behaviour can be changed:
\renewcommand\floatpagefraction{.9}
\renewcommand\topfraction{.9}
\renewcommand\bottomfraction{.9}
\renewcommand\textfraction{.1}   
\setcounter{totalnumber}{50}
\setcounter{topnumber}{50}
\setcounter{bottomnumber}{50}

\begin{document}

\begin{table}[H]
	\centering
	\begin{tabular}{|l||l|l||l|l|l|l|l|}
		\hline
		& \multicolumn{2}{c||}{\textbf{1-Ring $\mu$}} & \multicolumn{4}{c|}{\textbf{1-Ring $e$}} \\ \hline
		Error source & FHC & RHC & FHC & RHC & FHC 1 d.e. & FHC/RHC \\ \hline\hline
		SK Detector & 1.86 & 1.51 & 3.03 & 4.22 & 16.69 & 1.60 \\ \hline
		SK FSI+SI+PN & 2.20 & 1.98 & 3.01 & 2.31 & 11.43 & 1.57 \\ \hline
		Flux + Xsec constrained & 3.22 & 2.72 & 3.22 & 2.88 & 4.05 & 2.50 \\ \hline
		$\sigma(\nu_e)/\sigma(\bar{\nu}_e)$ & 0.00 & 0.00 & 2.63 & 1.46 & 2.62 & 3.03 \\ \hline
		NC1$\gamma$ & 0.00 & 0.00 & 1.08 & 2.59 & 0.33 & 1.49 \\ \hline
		NC Other & 0.25 & 0.25 & 0.14 & 0.33 & 0.98 & 0.18 \\ \hline
		Osc & 0.04 & 0.03 & 3.86 & 3.60 & 3.78 & 0.79 \\ \hline\hline
		All Systematics & 4.40 & 3.76 & 6.10 & 6.51 & 20.94 & 4.77 \\ \hline
		All with osc & 4.40 & 3.76 & 7.27 & 7.44 & 21.24 & 4.85 \\ \hline
	\end{tabular}
\end{table}


\end{document}
